\chapter{Conclusiones}
\label{cap:capitulo5}

\begin{flushright}
\begin{minipage}[]{10cm}
\emph{El éxito es la suma de pequeños esfuerzos repetidos día tras día.}\\
\end{minipage}\\

Robert Collier\\
\end{flushright}

\vspace{1cm}


En este último capítulo se exponen los objetivos alcanzados y las conclusiones obtenidas, así como las competencias y requisitos adquiridos durante el desarrollo del proyecto.
Se incluye también posibles líneas futuras de este trabajo.

\section{Objetivos y requisitos cumplidos}

En las siguientes subseciones se describirán los objetivos y requisitos cumplidos durante el desarrollo del presente proyecto.

\subsection{Objetivos}


Tras la realización de este proyecto se ha conseguido cumplir con el objetivo principal, que consistía en diseñar un sistema asistencial inteligente para persona mayores. Asimismo, se han alcanzado los subobjetivos planteados en la Sección~\ref{sec:descripcion}:

\begin{enumerate}
	\item Se han analizado las soluciones existentes en el ámbito de los robots asistenciales, utilizándolas de referencia para el diseño del sistema propuesto.
	\item Se han explorado distintas estrategias para desarrollar juegos interactivos, evaluando tanto plataformas especializadas como implementaciones propias.
	\item Se han estudiado y seleccionado los componentes \textit{hardware} y \textit{software} de menor coste sin perder en funcionalidad y compatibilidad.
	\item Se ha utilizado FreeCAD para el diseño y modelado de la estructura del robot.
	\item Se ha empleado una impresora 3D para la fabricación de las piezas físicas que forman la estructura del robot.
	\item Se ha utilizado un router propio para asegurar la conexión del sistema.
	\item Se han evaluado y comparado diferentes algoritmos de reconocimiento facial con el fin de conseguir la solución más adecuada para el sistema.
	\item Se han realizado pruebas comparativas  sobre distintos modelos de lenguaje para obtener la alternativa más adecuada adecuada para el sistema.
	\item Se han integrado los distintos módulos \textit{software} para conseguir un sistema funcional.
    \item Se han llevado a cabo experimentos en entornos reales para comprobar el correcto funcionamiento del sistema.
\end{enumerate}


\subsection{Requisitos}

A su vez, también se han cumplido los requisitos expresados en la Sección ~\ref{sec:Requisitos}:

\begin{enumerate}
	\item Mediante el uso de un ordenador externo como servidor, se han delegado los procesos pesados, permitiendo que la Raspberry Pi ejecute únicamente tareas de bajo coste computacional.
	\item Todas las librerías y herramientas \textit{software} empleadas son compatibles con la Raspberry Pi 4.
	\item Los componentes \textit{hardware} utilizados son accesibles y de bajo coste, manteniendo las funcionalidades requeridas para el proyecto.
	\item Las piezas físicas se han diseñado teniendo en cuenta las limitaciones de las impresoras 3D domésticas, facilitando su impresión en cualquier impresora estándar.
	\item La disposición triangular de las ruedas proporciona una base estable que permite el desplazamiento seguro del robot.
	\item El número y distribución de los componentes integrados en la estructura se han optimizado para mantener un peso equilibrado y una correcta distribución de la masa.
\end{enumerate}


\section{Competencias adquiridas}

Durante el desarrollo del trabajo se han adquirido conocimientos en diferentes áreas, entre las que destacan:

\begin{itemize}
	\item Profundización en el uso herramientas de diseño y modelado 3D mediante FreeCAD.
	\item Desarrollo de la capacidad de evaluar y seleccionar modelos y algoritmos en función de los requisitos establecidos.
	\item Aplicación de modularización del \textit{software} para su mantenimiento y depuración.
    \item Adquisición de conocimientos sobre arquitecturas clientes-servidor y el uso de API Getway para la comunicación entre servicios.
	\item Adquisición de conocimientos sobre el lenguaje de diseño CSS.
	\item Diseño y configuración de redes privadas locales para garantizar comunicaciones seguras entre dispositivos.
	\item Resolución de problemas de ensamblaje.
	\item Adquisición de conocimientos en electrónica, incluyendo la comprobación de tensiones, corrientes y conexión de  los componentes.
    \item Desarrollo de la capacidad de redactar un escrito en plataformas como \LaTeX.
\end{itemize}


\section{Conclusiones del desarrollado}

Durante el desarrollo de este proyecto, se han explorado diferentes áreas de conocimiento, como la comunicación entre dispositivos, la electrónica, el modelado 3D, entre otras. Gracias a ello, se han adquirido nuevos conocimientos y ampliando los previamente adquiridos.
\vspace{0.2cm}

Además, a lo largo del proyecto se ha desarrollado la capacidad para resolver problemas y adaptarse a los contratiempos que han surgido durante el desarrollo. Un ejemplo de ello fue la necesidad de emplear una arquitectura distribuida debido a las limitaciones computacionales de la Raspberry Pi, lo que permitió obtener un sistema funcional.

\vspace{0.2cm}

Asimismo, se ha concluido la posibilidad de desarrollar un sistema de interacción humano-robot fluido, natural y funcional mediante el empleo de componentes accesibles y de bajo coste.


\vspace{0.2cm}

Finalmente, se considera que soluciones como la desarrollada en este proyecto pueden complementar el apoyo y acompañamiento de personas mayores, favoreciendo su autonomía y mejorando su calidad de vida.


\section{Líneas futuras}

Para finalizar, se proponen diversas líneas de trabajo futuras que permitirían la continuidad de este sistema:

\begin{itemize}
	\item Incorporar un modelo de lenguaje comercial con mayor capacidad de razonamiento y compresión.
	\item Mejorar la interfaz gráfica para proporcionar al usuario una experiencia más intuitiva y personalizada.
	\item Utilizar un único núcleo de procesamiento, eliminando la dependencia de un servidor.
    \item Ampliar el sistema mediante la implementación de nuevas aplicaciones y funcionalidades.
    \item Implementar el seguimiento del usuario.
    \item Incorporar reconocimiento de emociones.
\end{itemize}

